\documentclass[11pt]{article}

\usepackage[a4paper,margin=25mm]{geometry}
\usepackage{array}
\usepackage{booktabs}
\usepackage{longtable}
\usepackage{tikzSteel}

\title{TikZ-Steel Manual}
\author{TikZ-Steel contributors}
\date{Development version}

\newcommand{\command}[1]{\texttt{\textbackslash #1}}

\begin{document}

\maketitle

\section{Overview}

TikZ-Steel is a TikZ-based package for drawing structural steel cross-sections.
The package currently focuses on cold-formed steel (CFS) and hot-rolled steel
(HRS) sections. It provides legacy positional commands, explicit CFS/HRS
positional wrappers, and a newer key-value interface.

\section{Loading the Package}

\begin{verbatim}
\usepackage{tikzSteel}
\end{verbatim}

The package loads TikZ and the TikZ \texttt{calc} library.

\section{Recommended Key-Value Interface}

The recommended public API uses commands named \command{TikZSteel...} with
readable keys:

\begin{verbatim}
\begin{tikzpicture}
  \TikZSteelLippedChannel[
    depth=245,
    flange=75,
    lip=20,
    thickness=2.5,
    radius=5,
    centerline=true,
    dimensions=true,
    label=C245,
    label x=35,
    label y=130,
    scale=0.2
  ]
\end{tikzpicture}
\end{verbatim}

\subsection{Common Keys}

\begin{longtable}{p{0.36\linewidth}p{0.56\linewidth}}
\toprule
Key & Meaning \\
\midrule
\texttt{depth} & Section depth or vertical size \\
\texttt{width} & Section width or horizontal size \\
\texttt{diameter} & Circular RC section diameter \\
\texttt{flange} & Common flange width \\
\texttt{top flange}, \texttt{bottom flange} & Asymmetric flange widths \\
\texttt{left flange}, \texttt{right flange} & Hat-section flange widths \\
\texttt{lip}, \texttt{top lip}, \texttt{bottom lip} & Lip lengths \\
\texttt{thickness} & Section thickness or line width in millimetres \\
\texttt{radius}, \texttt{root radius} & Bend or root radius \\
\texttt{gap} & Spacing between built-up members \\
\texttt{offset} & Offset between nested members \\
\texttt{position}, \texttt{stiffener position} & Relative stiffener position \\
\texttt{leg}, \texttt{vertical leg}, \texttt{horizontal leg} & Angle leg dimensions \\
\texttt{web depth}, \texttt{web thickness} & Welded section web dimensions \\
\texttt{flange thickness} & Rolled flange thickness \\
\texttt{top flange width}, \texttt{bottom flange width} & Welded section flange widths \\
\texttt{top flange thickness}, \texttt{bottom flange thickness} & Welded section flange thicknesses \\
\texttt{reference line width}, \texttt{line width} & Plate/bar reference line settings \\
\texttt{cover}, \texttt{bar diameter} & RC cover and reinforcement bar size \\
\texttt{top bars}, \texttt{bottom bars}, \texttt{left bars}, \texttt{right bars} & Rectangular RC bar counts \\
\texttt{top layers}, \texttt{bottom layers}, \texttt{side layers} & Rectangular RC reinforcement layers \\
\texttt{layer spacing}, \texttt{perimeter bars} & RC layer spacing and circular bar count \\
\texttt{label}, \texttt{label x}, \texttt{label y} & Optional label and label position \\
\texttt{filled}, \texttt{centerline}, \texttt{dimensions}, \texttt{monochrome}, \texttt{tie} & Optional drawing overlays and modes \\
\texttt{mode}, \texttt{simplified}, \texttt{detailed} & Accepted drawing-mode keys for future detail levels \\
\texttt{scale}, \texttt{rotate} & Local drawing transform \\
\texttt{x}, \texttt{y} & Local drawing offset \\
\bottomrule
\end{longtable}

\subsection{Current Dimension Overlay}

The \texttt{dimensions=true} key enables a basic visual overlay that uses the
\texttt{tikzSteel/dimension} style. In the current development version this
overlay is only a lightweight guide for showing generic horizontal and vertical
dimension lines around a section.

It is not yet a complete engineering dimensioning system. Dimension values,
extension lines, standard offsets, unit formatting, and standard-specific
notation are planned for later development. For now, the example galleries show
the command parameters beside each sketch as the clearer reference format.

\section{Cold-Formed Steel Commands}

\begin{longtable}{ll}
\toprule
Command & Purpose \\
\midrule
\command{TikZSteelChannel} & Unlipped channel \\
\command{TikZSteelLippedChannel} & Lipped channel \\
\command{TikZSteelStiffenedChannel} & Channel with an intermediate stiffener \\
\command{TikZSteelZee} & Zee section \\
\command{TikZSteelLippedZee} & Lipped zee section \\
\command{TikZSteelSigma} & Sigma-style stiffened channel \\
\command{TikZSteelHat} & Hat section \\
\command{TikZSteelAngle} & Angle section \\
\command{TikZSteelLippedAngle} & Lipped angle section \\
\command{TikZSteelRHS} & Rectangular hollow section \\
\command{TikZSteelSHS} & Square hollow section \\
\command{TikZSteelCHS} & Circular hollow section \\
\command{TikZSteelBackToBackChannels} & Built-up back-to-back channels \\
\command{TikZSteelToeToToeChannels} & Built-up toe-to-toe channels \\
\command{TikZSteelNestedChannels} & Nested channels \\
\command{TikZSteelBoxedChannels} & Boxed channels \\
\command{TikZSteelBuiltUpIChannels} & Built-up I from channels \\
\command{TikZSteelBackToBackZees} & Built-up back-to-back zees \\
\bottomrule
\end{longtable}

\section{Hot-Rolled Steel Commands}

\begin{longtable}{ll}
\toprule
Command & Purpose \\
\midrule
\command{TikZSteelUniversalBeam} & Universal beam/I-section \\
\command{TikZSteelUniversalColumn} & Universal column/I-section \\
\command{TikZSteelWeldedI} & Welded I-section \\
\command{TikZSteelTee} & Tee section \\
\command{TikZSteelHRSChannel} & Hot-rolled channel \\
\command{TikZSteelEqualAngle} & Equal angle \\
\command{TikZSteelUnequalAngle} & Unequal angle \\
\command{TikZSteelDoubleAngles} & Double angles \\
\command{TikZSteelPlate} & Plate or flat rectangle \\
\command{TikZSteelFlatBar} & Flat bar \\
\command{TikZSteelRoundBar} & Round bar \\
\command{TikZSteelHRSCHS} & Hot-rolled circular hollow section \\
\command{TikZSteelHRSRHS} & Hot-rolled rectangular hollow section \\
\command{TikZSteelHRSSHS} & Hot-rolled square hollow section \\
\bottomrule
\end{longtable}

\section{Reinforced Concrete Commands}

The initial RC commands draw simple rectangular and circular concrete
cross-sections with longitudinal rebars shown as filled circles. In this
development version, \texttt{cover} is interpreted as the sketch distance from
the concrete face to the rebar centerline or tie centerline.

\begin{longtable}{ll}
\toprule
Command & Purpose \\
\midrule
\command{TikZRCRectangular} & Rectangular RC section with top, bottom, and side bars \\
\command{TikZRCCircular} & Circular RC section with perimeter bars \\
\bottomrule
\end{longtable}

\begin{verbatim}
\TikZRCRectangular[
  width=300,
  depth=500,
  cover=40,
  bar diameter=18,
  top bars=2,
  bottom bars=3,
  bottom layers=2,
  layer spacing=28
]

\TikZRCCircular[
  diameter=450,
  cover=45,
  bar diameter=20,
  perimeter bars=10
]
\end{verbatim}

\section{Style Customization}

The default styles use red for straight segments and blue for curved segments.
They can be changed with \command{TikZSteelSetup}:

\begin{verbatim}
\TikZSteelSetup{
  tikzSteel/straight/.style={tikzSteel/default, black},
  tikzSteel/round/.style={tikzSteel/default, gray}
}
\end{verbatim}

\section{Examples}

\begin{center}
\begin{tikzpicture}
  \TikZSteelLippedChannel[
    depth=180,
    flange=55,
    lip=18,
    thickness=2,
    radius=4,
    centerline=true,
    dimensions=true,
    label=CFS,
    label x=30,
    label y=95,
    scale=0.04
  ]
  \TikZSteelUniversalBeam[
    depth=61.2,
    width=22.9,
    flange thickness=1.96,
    web thickness=1.19,
    root radius=1.4,
    x=9,
    scale=0.18
  ]
\end{tikzpicture}
\end{center}

\section{Development Status}

This manual is an initial development document. The next documentation milestone
is a full visual gallery with rendered examples, argument tables for every
legacy and key-value command, and a stable release-oriented API reference.

\end{document}
